% Depends on the full two-pole classification and its audited local formulas.
\section{The sum bound when one terminal pair is full}
\label{sec:one-full-tripole-sum}

\begin{lemma}[One full pair]
\label{lem:one-full-tripole-sum}
Let $T$ be a connected loopless multigraph of rank $h$.
Its three distinct terminals $a,b,c$ have degree two.
Every other vertex has degree three.
If $P_{ac}(T)=2^h$, then
\[
 P_{ab}(T)+P_{bc}(T)\le2^h+6\,2^{h/2}.
\]
\end{lemma}

\begin{proof}
Suppress $b$ to an edge of a two-terminal network $N$.
Its rank remains $h$, and $P_{ac}(N)=2^h$.
The graph $N$ has no loop, since every terminal path would avoid that loop.
Lemma~\ref{lem:full-two-pole-classification} therefore applies.
Its positive-rank bridgeless pieces form a series chain.
Each piece is a capped ladder with singleton and square blocks.

First suppose $N$ consists of one such piece.
For a prefix with $p$ ladder positions, write $A_p$ for its cap count.
The empty prefix has $A_0=1$.
Adding a singleton replaces $A$ by $A+1$.
Adding a square replaces $A$ by $2A+2$.
These recurrences give
\begin{equation}\label{eq:tripole-prefix-cap}
 A_p\le3\,2^{p/2}-2\qquad(p\ge0).
\end{equation}
The singleton step uses
$3(\sqrt2-1)2^{p/2}\ge1$.
The square step preserves equality in the bound.
The same estimate holds for a suffix cap count $B_q$.

Let $p$ and $q$ count positions strictly before and after the block containing
the subdivided edge. For an interface, they count the positions on its two sides.
The subdivision formulas in Lemma~\ref{lem:one-full-three-pole-gap} give
the following table:
\[
\begin{array}{c|c|c|c}
\text{edge type}&h&P_{bc}&P_{ab}\\ \hline
\text{interface}&p+q+1&(1+A_p)2^q&(1+B_q)2^p\\
\text{singleton rung}&p+q+2&2(1+A_p)2^q&2(1+B_q)2^p\\
\text{square edge}&p+q+3&(4+3A_p)2^q&(4+3B_q)2^p
\end{array}
\]
The two parallel edges form the rank-one empty-ladder case.
Subdividing one gives $P_{ab}=P_{bc}=2$, which satisfies the conclusion.

For the first two rows, divide the path counts by $2^h$.
Their sum is bounded by $f_p+f_q$, where
\[
 f_j=\frac{1+A_j}{2^{j+1}}.
\]
Here the suffix uses its own cap count in the second summand.
Equation~\eqref{eq:tripole-prefix-cap} gives $f_j\le1/2$ for $j\ge3$.
Thus the required sum is at most $2^h$ if both $p,q\ge3$.
Otherwise reverse the piece if necessary and assume $p\le2$.
Bound $P_{bc}$ by $2^h$ using the affine path count.
In the interface row, the other term is at most
\[
 3\,2^{p+q/2}=3\,2^{(p-1)/2}2^{h/2}
 \le3\sqrt2\,2^{h/2}.
\]
In the singleton row, it is at most
\[
 6\,2^{p+q/2}=6\,2^{(p-2)/2}2^{h/2}
 \le6\,2^{h/2}.
\]

For a square edge, the normalized counts use
\[
 g_j=\frac{4+3A_j}{2^{j+3}}.
\]
Equation~\eqref{eq:tripole-prefix-cap} gives $g_j\le1/2$ for $j\ge2$.
If both $p,q\ge2$, their sum is at most one.
Otherwise assume $p\le1$ and again bound $P_{bc}$ by $2^h$.
Since $4+3B_q\le9\,2^{q/2}$, the remaining term is at most
\[
 9\,2^{p+q/2}=9\,2^{(p-3)/2}2^{h/2}
 \le\frac92\,2^{h/2}.
\]
This proves the bound for one piece.

Now suppose the subdivided edge lies inside a piece of rank $t$.
Let $L$ and $R$ be the sums of the piece ranks before and after this piece.
Write $A$ and $B$ for its paths from $b$ to its left and right endpoints.
The series decomposition gives
\[
 h=L+t+R,\qquad P_{ab}=2^L A,\qquad P_{bc}=2^R B.
\]
If $L,R\ge1$, each count is at most $2^{h-1}$.
If exactly one side has positive rank, reverse the chain and write $L=0$.
The one-piece result and $B\le2^t$ give
\[
 A+2^R B
 =(A+B)+(2^R-1)B
 \le2^{t+R}+6\,2^{t/2}
 \le2^h+6\,2^{h/2}.
\]
The case $L=R=0$ was already proved.

Finally, suppose the subdivided edge is a bridge of $N$.
Its two sides have positive ranks $L,R$ with $h=L+R$.
All terminal paths on each side are full, so
\[
 P_{ab}+P_{bc}=2^L+2^R\le2^{L+R}=2^h.
\]
This completes the proof.
\end{proof}

\begin{lemma}[One full pair leaves a small adjacent pair]
\label{lem:one-full-tripole-small-adjacent}
Under the hypotheses of Lemma~\ref{lem:one-full-tripole-sum},
\[
 \min\{P_{ab},P_{bc}\}\le3\,2^{3h/4}.
\]
\end{lemma}

\begin{proof}
Use the same full two-pole decomposition and local table.
Suppose first that the subdivided edge lies in a nonempty capped piece.
Its rank is $t=p+q+c$, where $c=1,2,3$ for the three table rows.
Equation~\eqref{eq:tripole-prefix-cap} bounds the two normalized local counts by
\[
 \frac32\,2^{-q/2},\qquad \frac32\,2^{-p/2},
\]
respectively. The square row actually has the smaller coefficient $9/8$.
Let $L,R$ be the sums of the ranks before and after this piece.
The two global normalized counts are therefore at most
\[
 \frac32\,2^{-R-q/2},\qquad \frac32\,2^{-L-p/2}.
\]
The sum of these two exponents is
\[
 L+R+\frac{p+q}{2}
 =h-\frac t2-\frac c2\ge\frac h2-\frac c2.
\]
At least one exponent is at least $h/4-c/4$.
Since $(3/2)2^{3/4}<3$, the smaller normalized count is at most
$3\,2^{-h/4}$.

If the active piece is a digon, both local counts are full.
The two global normalized counts are $2^{-R},2^{-L}$, with $L+R=h-1$.
Their minimum is at most $2^{-(h-1)/2}\le3\,2^{-h/4}$.
If the subdivided edge is a bridge, the same argument uses $L+R=h$.
Multiplication by $2^h$ proves the result.
\end{proof}
